\begin{exercise}{147}
Let $f\colon \mathcal{C}\to \mathcal{D}$ and $g\colon \mathcal{D}\to \mathcal{C}$ be functors. Suppose that $\varepsilon\colon fg\to\mathrm{id}$ is a
natural transformation. Show that the map
\[
 \mathrm{map}(x,g(y))\to\mathrm{map}(f(x),f(g(y)))\to\mathrm{map}(f(x),y)
\]
is natural in $x$.

\begin{proof}
Fix $y\in \mathcal{D}$. By Definition~4.2.3, $f$ induces a morphism of right
fibrations over $\mathcal{C}\times \mathcal{C}^{\mathrm{op}}$,
\[
 \mathrm{Tw}(\mathcal{C})\longrightarrow
 (\mathcal{C}\times \mathcal{C}^{\mathrm{op}})
 \times_{\mathcal{D}\times \mathcal{D}^{\mathrm{op}}}\mathrm{Tw}(\mathcal{D}).
\]
Here the map of bases is $f\times f^{\mathrm{op}}$.
Under the dual of Theorem~3.3.16, this is a natural transformation
\[
 \mathrm{map}_{\mathcal{C}}(-,-)\longrightarrow\mathrm{map}_{\mathcal{D}}(f(-),f(-))
\]
of functors $\mathcal{C}^{\mathrm{op}}\times \mathcal{C}\to\mathrm{Spc}$
(Definition~4.2.5). Its component at $(x,z)$ sends an arrow
$a\colon x\to z$ to $f(a)$, since the map on twisted arrows is
postcomposition with $f$ in every simplicial degree.
Restricting to $z=g(y)$ gives the first arrow of the required
natural transformation.

For the second arrow, apply the mapping bifunctor of $\mathcal{D}$ to
\[
 \mathcal{C}^{\mathrm{op}}\times\Delta^1
 \xrightarrow{\ f^{\mathrm{op}}\times\varepsilon_y\ }
 \mathcal{D}^{\mathrm{op}}\times \mathcal{D}.
\]
\begingroup
\setlength{\emergencystretch}{2em}
This gives a natural transformation from
$\mathrm{map}_{\mathcal{D}}(f(-),f(g(y)))$ to $\mathrm{map}_{\mathcal{D}}(f(-),y)$,
whose components are postcomposition with $\varepsilon_y$.
Composing these two transformations in
$\mathrm{Fun}(\mathcal{C}^{\mathrm{op}},\mathrm{Spc})$ gives exactly the displayed
composite, with its full naturality in $x$.
\par
\endgroup
\end{proof}
\end{exercise}
